% Circonferenza goniometrica di raggio 1: il punto P all'angolo theta ha
% coordinate (cos theta, sin theta). Disegnato theta = 120 gradi, con cos < 0,
% e tratteggiato il punto P' a 60 gradi (180 - theta), simmetrico rispetto a y.
\begin{tikzpicture}[>=Stealth, font=\small, scale=1.35]
  \draw[->] (-1.35,0) -- (1.45,0) node[right] {$x$};
  \draw[->] (0,-1.25) -- (0,1.35) node[above] {$y$};
  \draw[thick, richiamo] (0,0) circle (1);
  \node[below right] at (1,0) {$1$};
  \node[above left] at (0,1) {$1$};
  \coordinate (P) at (120:1);
  \coordinate (Q) at (60:1);
  \draw[dashed, formula] (0,0 |- P) -- (Q);
  \draw[formula, fill=white] (Q) circle (1.4pt) node[above right] {$P'$};
  \draw[very thick, formula] (0,0) -- (P);
  \fill[formula] (P) circle (1.2pt) node[above left] {$P$};
  \draw[dashed] (P) -- (P |- 0,0);
  \draw[dashed] (P) -- (0,0 |- P);
  \draw[esempio, thick] (0.3,0) arc[start angle=0, end angle=120, radius=0.3];
  \node[esempio] at (35:0.45) {$\theta$};
  \draw[very thick, medica] (0,0) -- (P |- 0,0);
  \node[medica, below] at (-0.62,0) {$\cos\theta$};
  \draw[very thick, attenzione] (0,0) -- (0,0 |- P);
  \node[attenzione, right] at (0,0.55) {$\sin\theta$};
\end{tikzpicture}
